% ========================================================================================================================================================================================
\section{Results}
This work produced two sets of results: firstly, those derived from
 optimisation conducted via sampling directly from the latent objective
 function described by the physical system and secondly, those stemming
from the optimisation of the five surrogate functions trained to emulate 
 the latent objective function and the associated physical system.
The latter set of results were generated to provide statistical support for
 the former.

% ========================================================================================================================================================================================
\subsection{Physical System Optimisation}
A single Bayesian sequential optimisation experiment
 was run, with a sample limit (\(p_{\text{tot}}=27\)), number of samples
 drawn per sequential iteration (\(p_{\text{iter}}=3\)), and a prior-forming
 dataset of 8 quasirandomly sampled targets 
 (Fig.~\ref{fig:BayesianPriorsAbbreviated}).
This Bayesian routine employed the probability of improvement acquisition
 function (Eq.~\ref{eq:ProbabilityOfImprovement}).
Figure \ref{fig:ResultOverview} illustrates the progression of this
 sequential routine when compared with the one-shot pseudorandom
 sampling technique, with the mass losses plotted on the y-axis being
 the best values discovered given the samples evaluated so far
 (\(\Delta m_{\%}^{\circ}\)).
A simple statistical review of the experimental datasets generated by 
 both the sequential Bayesian and one-shot pseudorandom techniques within
 the sample limit is provided in Table~\ref{tab:BayesianSequentialsBestValues}.
Given the sequential samples available (\(p_{\text{tot}}-p_{\text{init}}=19\))
 and the samples per iteration (\(p_{\text{iter}}=3\)), an extra iteration of 
 sequential sampling was run which consequently breached the sample limit-
 these extra samples are included in Figure \ref{fig:ResultOverview}, but
 are not considered in the preparation of Table 
 \ref{tab:BayesianSequentialsBestValues} (as it would be unfair to compare
 two different sampling techniques which were allowed an unequal number of
 samples over which to optimise).

\begin{center}
    \captionsetup{format=custom,width=0.48\textwidth,skip=0pt}
    \includegraphics[width=0.48\textwidth]{PracticalResultOverview.png}
    \captionof{figure}{
        \footnotesize{
            Plot of optimisation progression for the experimentally
             deployed Bayesian optimisation routine utilising the
             probability of improvement (PI, solid blue line) 
             acquisition function, and compared alongside the 
             one-shot pseudorandom sampling technique 
             (PR, dashed line).
            Uncertainty in \(\Delta\text{m}_{\text{\%}}\) values
             is equal to 0.44\% points (\(1\sigma\)).
        }
    }
    \label{fig:ResultOverview}
\end{center}

Figure \ref{fig:ResultOverview} shows that although the sequential 
 Bayesian technique discovers a sample surpassing the pseudorandom
 benchmark by sample 21 (77\% of \(p_{\text{tot}}\)), its first 
 statistically significant improvement was recorded on sample 27
 (100\% of \(p_{\text{tot}}\)) - this was consistent with 
 exploratory-biased behaviour.
Further evidence of the probability of improvement acquisition
 functions exploratory-biased behaviour is seen in 
 Table~\ref{tab:BayesianSequentialsBestValues}, where, despite an
 improvement upon the initial prior's highest discovered target value
 (14.01\% points \(\rightarrow\) 16.67\% points), the running mean of
 samples within the growing dataset fell from its quasirandomly
 sampled prior's high (12.82\% points) to its final lowpoint (12.50\% 
 points) which reached an approximate parity with the pseudorandomly
 sampled benchmark dataset (12.54\% points).
Contrary to our initial expectations, the evidence presented above shows 
 that the probability of improvement acquisition function does not 
 exhibit overly exploitative behaviour due to the lack of a marginal 
 threshold parameter \(\varepsilon\) being manually tuned.

\footnotesize
\begin{center}
    \begin{tblr}{colspec={Q[si={table-format=3.2},l]Q[si={table-format=3.2},l]Q[si={table-format=3.2},l]Q[si={table-format=3.2},l]Q[si={table-format=3.2},l]},row{1}={guard},column{1}={guard}}
        {Technique\\Function}&{No.}&{Max\\(\%)}&{Mean\\(\%)}&{Standard\\Dev. (\%)}\\
        \hline
        P  &   8     & 14.01 & 12.82 & 1.04 \\
        \hline[dashed]
        PI &   27     & 16.67 & 12.50 & 1.63 \\
        PR &   27     & 14.01 & 12.54 & 0.88 \\
    \end{tblr}
    \captionsetup{format=custom,width=0.48\textwidth,skip=0pt}
    \captionof{table}{
        \footnotesize{
            Statistical review of datasets generated experimentally
             within the sample limit.
            P: Prior dataset, quasirandomly sampled and 
             fed into Bayesian routine.
            PI: Final dataset from Bayesian routine
             using probability of improvement acquisition function.
            PR: Dataset from one-shot pseudorandom sampler.
            No.: Number of samples in dataset.
            Max: Highest discovered mass loss (\(\Delta m_{\%}^{\circ}\))
            Mean: Population mean.
            Standard Dev.: A single standard deviation.
        }
    }
    \label{tab:BayesianSequentialsBestValues}
\end{center}
\normalsize

Notably, the Bayesian sequential technique employing the probability of
 improvement acquisition function discovered samples attaining the highest
 mass losses of any sample characterised during this work
 (\(\Delta m_{\%}^{\circ}=20.77\%\)), albeit outside the sample limit
 (See Figure \ref{fig:ResultOverview}).

% ========================================================================================================================================================================================
\subsection{Emulated System Optimisation}
Since curing of a single prepolymer sample could take up to 
 a week to complete, the temporal expense of sampling the objective function was
 too high to allow for repeat experiments involving the sequential
 techniques deployed.
This precluded statistically proving the utility of Bayesian optimisation
 in this field experimentally within reasonable 
 timeframes.
However, by developing inexpensive surrogate functions capable 
 of emulating the physical system, a much larger number of function 
 evaluations could be carried out computationally, facilitating statistical 
 analysis of sampling methods' efficacy.
By despatching 1,000 simulations of any given sampling technique's ability
 to optimise any given emulator's surrogate function and returning an
 array of maximum target variable values \(\Delta m_{\%}^{\circ}\)
 discovered, basic statistical conclusions could be drawn.
All simulated sampling techniques observed a sample limit of \(p_{tot}=27\),
 whilst the sequential techniques were all initiated using \(p_{init}=8\)
 quasirandomly drawn samples.
Sampling techniques spanned one-shot and sequential forms, with the
 former group comprised of pseudorandom (PR) and quasirandom sampling (QR)
 whilst the latter was made up of both objective-naive and Bayesian methods.
The lone objective-naive space-filling technique simulated was the 
 mc-intersite-proj-th algorithm (MIPT), whilst the two Bayesian techniques 
 diverged in terms of the covariance function employed - either probability 
 of improvement (PI) or expected improvement (EI).

\begin{figure*}[htb]
    \captionsetup{format=custom,width=0.8\textwidth,skip=0pt}
    \includegraphics[width=\textwidth]{ComputationalResultOverview.png}
    \captionof{figure}{
        \footnotesize{
            Strip plots of maximum function evaluations obtained
             \(\Delta m_{\%}^{\circ}\), for each of the 25,000
             simulated sampling exercises undertaken.
            Every simulated sampling exercise observed a constant
             sample limit \(p_{tot}=27\), with each sampling
             technique being simulated 1,000 times upon each of
             the emulators, and the Bayesian routines were initiated
             with 8 quasirandomly drawn samples.
        }
    }
    \label{fig:CompResOverv}
\end{figure*}

% \begin{figure*}[htb]
%     \captionsetup{format=custom,width=0.8\textwidth,skip=0pt}
%     \includegraphics[width=\textwidth]{SimulationResultOverview}
%     \captionof{figure}{
%         \footnotesize{
%             Swarm plot of maximum function evaluations obtained,
%             \(\Delta m_{\%}^{\circ}\), for each of the 4,000 simulated 
%             optimisations carried out.
%             All optimisations used 27 samples, each technique was 
%             simulated 1,000 times, and the Bayesian routines were
%             initiated with 8 quasirandomly sampled points.
%         }
%     }
%     \label{fig:SimulationResultOverview}
% \end{figure*}

Figure~\ref{fig:CompResOverv} presents a series of strip plots showing
 distribution of outcomes for the 25,000 simulations undertaken across
 the five sampling techniques and five emulators studied.
On comparison of the mean \(\Delta m_{\%}^{\circ}\) values discovered by the
 various sampling techniques, sequential Bayesian techniques were shown to
 generally perform better (as observed in 4/5 emulators trialled) than either
 the one-shot or objective naive sequential techniques.
Although standard deviation values associated with these mean values would
 preclude a straightforward conclusion that the difference between values
 are statistically significant, it can be seen that the simulations are 
 distributed in a somewhat lognormal fashion which shows that only the 
 Bayesian techniques are capable of discovering well-above-average-performing 
 samples (causing the majority of skew seen in the means).

To investigate the degree to which exceptionally high value formulations
 were uncovered by the various sampling techniques, a `region of global 
 optimality' was devised for each of the emulators, with an upper bound 
 defined by the computational global maximum and a lower bound defined as 
 being 10\% lower than this (as relativised by the computational global 
 minimum).
The Bayesian sequential technique adopting the expected improvement
 acquisition function in its sampling policy not only outperformed the
 probability of improvement technique in terms of mean values attained,
 but also in number of simulations obtaining \(\Delta m_{\%}^{\circ}\)
 values which resided within the small region of global optimality,
 as evidenced in Table~\ref{tab:CompResOvervTab}.
Notably, neither of the one-shot techniques nor the objective-naive
 sequential technique were able to produce as high a proportion of 
 samples within the `region of global optimality', with all three
 consistently producing either one or two orders of magnitude less
 samples than that generated by the Bayesian techniques 
 (Table~\ref{tab:CompResOvervTab}).

Taking the GP-Mat-05 as the most accurate of the 
 emulators, and subject to the assumptions that it adequately represents
 the underlying objective function and that the defined `region of global
 optimality' is a reasonable proxy for the experimentally identified 
 optimum, the estimated discovery probabilities were 65\% for expected
 improvement and 23.8\% for probability of improvement.
This would make the chances of having discovered these
 exceptionally well-performing formulations during
 the course of our one-off implementation of the Bayesian sequential
 routine significantly more likely than the 0.6\%, 1.3\%, and 0.5\%
 probabilities exhibited by the two one-shot techniques (PR and QR) or the
 objective-naive MIPT technique respectively.

\footnotesize
\begin{center}
    \begin{tblr}{colspec={Q[si={table-format=3.2},l]Q[si={table-format=3.2},c]Q[si={table-format=3.2},c]Q[si={table-format=3.2},c]Q[si={table-format=3.2},c]Q[si={table-format=3.2},c]},row{1}={guard},column{1}={guard}}
        {Emu.}&{PR \\ (\%)}&{QR \\ (\%)}&{MIPT \\ (\%)}&{PI \\ (\%)}&{EI \\ (\%)}\\
        \hline
        M05    & 0.6    & 1.3       & 0.5   & 23.8  & 65.0 \\
        M15    & 6.9    & 5.9       & 3.4   & 26.3  & 69.9 \\
        M25    & 9.1    & 6.9       & 5.3   & 26.6  & 59.3 \\
        Rbf    & 65.2   & 71.1      & 53.0  & 96.8  & 100.0 \\
        Xgb    & 0.8    & 0.4       & 0.4   & 1.3   & 5.1 \\
    \end{tblr}
    \captionsetup{format=custom,width=0.48\textwidth,skip=0pt}
    \captionof{table}{
        \footnotesize{
            Summary of percent of simulations surpassing the
             lower region of optimality bound.
            Emu.: Emulators simulated upon.
            MO5/M15/M25: Gaussian process emulators
             where \(\nu=0.5/1.5/2.5\) respectively.
            Rbf: Gaussian process emulator with RBF kernel.
            Xgb: XGBoost gradient boosted decision tree emulator.
        }
    }
    \label{tab:CompResOvervTab}
\end{center}
\normalsize
